% GENERATED FILE. Edit canonical lessons, then run npm run generate:books.
% canonical-source-hash: fnv1a64:bf6c9c5a1290252b
% canonical-lessons: RU-C419-registratsiya, RU-C419-registrirovatsya, RU-C419-reys, RU-C419-kladi, RU-C419-posadochnyy

\chapter{Checking in for a flight}
\label{ch:ru-checking-in-for-a-flight}
\hlchaptermodality{419}

\begin{chapteropening}
\textbf{By the end of this chapter:} \emph{I can check in for a flight: check-in, checking in, the flight, hand luggage and the boarding pass.}

Complete the last lesson of chapter 419: I can check in for a flight: check-in, checking in, the flight, hand luggage and the boarding pass.
\end{chapteropening}

\section[\texorpdfstring{\ru{регистрация} (registrátsiya)}{регистрация (registrátsiya)}]{\ru{регистрация} (registrátsiya) — check-in; registration}

\label{lesson:RU-C419-registratsiya}



\begin{quote}
Before the new one: say the Russian for to go past, then the Russian for to see off.
\end{quote}

\subsection*{You'll want to know: \ru{регистрация}}
\textbf{\ru{регистрация}} (\emph{registrátsiya}) — \textquotedblleft{}check-in; registration\textquotedblright{}.

\begin{quote}
\ru{Регистрация} \ru{начинается} \ru{в} \ru{шесть} \ru{утра}.
\end{quote}

\begin{quote}
\emph{Registrátsiya nachináyetsya v shest utrá.}
\end{quote}

\begin{quote}
Check-in opens at six in the morning.
\end{quote}

Feminine. \emph{\ru{Стойка} \ru{регистрации}} — the check-in desk.

\begin{cousinweb}
\emph{\ru{регистрация}} is a learned word built on \emph{\ru{регистр}} 'a register', usually traced to Medieval Latin \emph{registrum} 'a list'. EtymDB lists the Russian word without a source, so the line is hedged.
\end{cousinweb}

\begin{grammarlens}[title={What you've built}]
One more word for telling, explaining and describing at B1.
\end{grammarlens}

\subsection*{Your turn}
\begin{itemize}
\raggedright
  \item \emph{Say it:} \emph{registrátsiya}
  \item \emph{Say it:} \emph{registrátsiya}, once more
  \item \emph{Say it:} a sentence of your own with \emph{registrátsiya}
  \item \emph{Recall:} say the Russian for to go past, then the Russian for to see off, then say \emph{registrátsiya} again
\end{itemize}

\begin{tcolorbox}[breakable,colback=teal!4,colframe=teal!35!black,title={Before you move on}]
What does \ru{регистрация} mean? (\textquotedblleft{}Check-in; registration\textquotedblright{}.) Say it once more.
\end{tcolorbox}
\section[\texorpdfstring{\ru{регистрироваться} (registrírovatsya)}{регистрироваться (registrírovatsya)}]{\ru{регистрироваться} (registrírovatsya) — to check in; to register}

\label{lesson:RU-C419-registrirovatsya}



\begin{quote}
Before the new one: say the Russian for to see off, then the Russian for check-in.
\end{quote}

\subsection*{You'll want to know: \ru{регистрироваться}}
\textbf{\ru{регистрироваться}} (\emph{registrírovatsya}) — \textquotedblleft{}to check in; to register\textquotedblright{}.

\begin{quote}
\ru{Сейчас} \ru{все} \ru{регистрируются} \ru{на} \ru{сайте}.
\end{quote}

\begin{quote}
\emph{Seychás vse registríruyutsya na sáyte.}
\end{quote}

\begin{quote}
Nowadays everyone checks in on the website.
\end{quote}

Imperfective; present \emph{\ru{регистрируюсь}}. Takes \emph{\ru{на}} + accusative for a flight. The perfective is \emph{\ru{зарегистрироваться}}.

\begin{cousinweb}
\emph{\ru{регистрироваться}} is built on \emph{\ru{регистр}} 'a register' with \emph{-\ru{ировать}} and \emph{-\ru{ся}}. \emph{\ru{Регистр}} is usually traced to Medieval Latin \emph{registrum} 'a list'; EtymDB lists \emph{\ru{регистрировать}} without a source.
\end{cousinweb}

\begin{grammarlens}[title={What you've built}]
One more word for telling, explaining and describing at B1.
\end{grammarlens}

\subsection*{Your turn}
\begin{itemize}
\raggedright
  \item \emph{Say it:} \emph{registrírovatsya}
  \item \emph{Say it:} \emph{registrírovatsya}, once more
  \item \emph{Say it:} a sentence of your own with \emph{registrírovatsya}
  \item \emph{Recall:} say the Russian for to see off, then the Russian for check-in, then say \emph{registrírovatsya} again
\end{itemize}

\begin{tcolorbox}[breakable,colback=teal!4,colframe=teal!35!black,title={Before you move on}]
What does \ru{регистрироваться} mean? (\textquotedblleft{}To check in; to register\textquotedblright{}.) Say it once more.
\end{tcolorbox}
\section[\texorpdfstring{\ru{рейс} (reys)}{рейс (reys)}]{\ru{рейс} (reys) — a flight; a trip (of a bus or a ship)}

\label{lesson:RU-C419-reys}



\begin{quote}
Before the new one: say the Russian for check-in, then the Russian for to check in.
\end{quote}

\subsection*{You'll want to know: \ru{рейс}}
\textbf{\ru{рейс}} (\emph{reys}) — \textquotedblleft{}a flight; a trip (of a bus or a ship)\textquotedblright{}.

\begin{quote}
\ru{Наш} \ru{рейс} \ru{в} \ru{два} \ru{часа} \ru{дня}.
\end{quote}

\begin{quote}
\emph{Nash reys v dva chasá dnya.}
\end{quote}

\begin{quote}
Our flight is at two in the afternoon.
\end{quote}

Masculine. \emph{\ru{Рейс} \ru{отменили}} — the flight was cancelled.

\begin{cousinweb}
\emph{\ru{рейс}} was borrowed from Dutch \emph{reis} 'a journey, travel' (EtymDB); in Russian it names one trip of a plane, a ship or a bus.
\end{cousinweb}

\begin{grammarlens}[title={What you've built}]
One more word for telling, explaining and describing at B1.
\end{grammarlens}

\subsection*{Your turn}
\begin{itemize}
\raggedright
  \item \emph{Say it:} \emph{reys}
  \item \emph{Say it:} \emph{reys}, once more
  \item \emph{Say it:} a sentence of your own with \emph{reys}
  \item \emph{Recall:} say the Russian for check-in, then the Russian for to check in, then say \emph{reys} again
\end{itemize}

\begin{tcolorbox}[breakable,colback=teal!4,colframe=teal!35!black,title={Before you move on}]
What does \ru{рейс} mean? (\textquotedblleft{}A flight; a trip (of a bus or a ship)\textquotedblright{}.) Say it once more.
\end{tcolorbox}
\section[\texorpdfstring{\ru{кладь} (klad)}{кладь (klad)}]{\ru{кладь} (klad) — luggage (in \ru{ручная} \ru{кладь}: hand luggage)}

\label{lesson:RU-C419-kladi}



\begin{quote}
Before the new one: say the Russian for to check in, then the Russian for a flight.
\end{quote}

\subsection*{You'll want to know: \ru{кладь}}
\textbf{\ru{кладь}} (\emph{klad}) — \textquotedblleft{}luggage (in \ru{ручная} \ru{кладь}: hand luggage)\textquotedblright{}.

\begin{quote}
\ru{В} \ru{самолёт} \ru{можно} \ru{взять} \ru{одну} \ru{сумку} \ru{ручной} \ru{клади}.
\end{quote}

\begin{quote}
\emph{V samolyót mózhno vzyat odnú súmku ruchnóy kládi.}
\end{quote}

\begin{quote}
You can take one bag of hand luggage onto the plane.
\end{quote}

Feminine; formal, and mostly in \emph{\ru{ручная} \ru{кладь}} — hand luggage.

\begin{cousinweb}
\emph{\ru{кладь}} is built on the root of \emph{\ru{класть}} 'to put, to lay': the things one loads. EtymDB lists \emph{\ru{кладь}} without a deeper source.
\end{cousinweb}

\begin{grammarlens}[title={What you've built}]
One more word for telling, explaining and describing at B1.
\end{grammarlens}

\subsection*{Your turn}
\begin{itemize}
\raggedright
  \item \emph{Say it:} \emph{klad}
  \item \emph{Say it:} \emph{klad}, once more
  \item \emph{Say it:} a sentence of your own with \emph{klad}
  \item \emph{Recall:} say the Russian for to check in, then the Russian for a flight, then say \emph{klad} again
\end{itemize}

\begin{tcolorbox}[breakable,colback=teal!4,colframe=teal!35!black,title={Before you move on}]
What does \ru{кладь} mean? (\textquotedblleft{}Luggage (in \ru{ручная} \ru{кладь}: hand luggage)\textquotedblright{}.) Say it once more.
\end{tcolorbox}
\section[\texorpdfstring{\ru{посадочный} (posádochnyy)}{посадочный (posádochnyy)}]{\ru{посадочный} (posádochnyy) — boarding (adjective)}

\label{lesson:RU-C419-posadochnyy}



\begin{quote}
Before the new one: say the Russian for a flight, then the Russian for luggage.
\end{quote}

\subsection*{You'll want to know: \ru{посадочный}}
\textbf{\ru{посадочный}} (\emph{posádochnyy}) — \textquotedblleft{}boarding (adjective)\textquotedblright{}.

\begin{quote}
\ru{Без} \ru{посадочного} \ru{талона} \ru{в} \ru{самолёт} \ru{не} \ru{пустят}.
\end{quote}

\begin{quote}
\emph{Bez posádochnogo talóna v samolyót ne pústyat.}
\end{quote}

\begin{quote}
They will not let you onto the plane without a boarding pass.
\end{quote}

\emph{\ru{Посадочный} \ru{талон}} — a boarding pass. \emph{\ru{Посадка}} is boarding, and also landing.

\begin{cousinweb}
\emph{\ru{посадочный}} is built on \emph{\ru{посадка}} 'boarding; landing', from \emph{\ru{по}-} and \emph{\ru{садить}} 'to seat'. EtymDB lists \emph{\ru{садить}} without a deeper source.
\end{cousinweb}

\begin{grammarlens}[title={What you've built}]
One more word for telling, explaining and describing at B1.
\end{grammarlens}

\subsection*{Your turn}
\begin{itemize}
\raggedright
  \item \emph{Say it:} \emph{posádochnyy}
  \item \emph{Say it:} \emph{posádochnyy}, once more
  \item \emph{Say it:} a sentence of your own with \emph{posádochnyy}
  \item \emph{Recall:} say the Russian for a flight, then the Russian for luggage, then say \emph{posádochnyy} again
\end{itemize}

\begin{tcolorbox}[breakable,colback=teal!4,colframe=teal!35!black,title={Before you move on}]
What does \ru{посадочный} mean? (\textquotedblleft{}Boarding (adjective)\textquotedblright{}.) Say it once more.
\end{tcolorbox}
